% Track d3 -- The unity and its differentiation in motion: the closure of the force law.
% Verification: investigations/rest_energy/verify/d3_closure_dynamics.py (claims d3-0 ... d3-23);
% independent adversarial check: investigations/rest_energy/verify/x3_check.py (x3-0 ... x3-23).
\section{The differentiation in motion: rotor, rate and duality angle}
\label{d3:sec:main}

This section follows the undivided energy and its differentiation in motion inside the one dynamics that the arena closes by itself, the
Lie--Poisson closure of the force law (Section~17 of \cite{ES17}), in which the rate $\kappa=c/\ell$ of $\alpha_\pm=\partial_\tau\pm\kappa G$, where the
hypothesis puts $c^2$, is the angular velocity of a rotor. No Yang--Mills field is used. As in \cref{d2:prop:duality}, two fields are kept apart: the
\emph{momentum field} $\pi=\eta^{-1}p$ ($\pi=m\alpha_+X$ in the clamped interior) and the \emph{frame field} $X$. In brief: the closure freezes the
undivided energy, rest energy and Poynting vector of the momentum field, so no Hamiltonian generates or changes the rest energy of that field,
while the rest energy of the frame field oscillates in the clamped interior and grows without bound under the free-rotor drift
(\cref{d3:thm:frozen,d3:cor:nogen}). The conserved total splits into a rate part $\tfrac12(I_0/\ell^2)c^2$ and a Lorentz-invariant frame part with
an amplitude set by the rest energy of the momentum field, and the duality angle swings with $\Omega^2$ proportional to it
(\cref{d3:prop:pendulum,d3:prop:omega,d3:cor:flat,d3:prop:budget}, restating Theorem~17.13 and Corollaries~17.14--17.15 of \cite{ES17}). A constant
rate is the heavy-rotor limit (\cref{d3:prop:multiscale}); nothing in the algebra fixes $\ell$ (\cref{d3:thm:noscale}); and $\hbar\kappa=mc^2$,
equivalently $\ell=\hbar/(mc)$, is an added identification, not a result, under which a quantised rotor has the first level $\tfrac12mc^2$
(\cref{d3:prop:norel,d3:prop:ladder}).

\paragraph{Conventions.} As in Section~17 of \cite{ES17}: $V=\R^{3,3}$, $\eta=\operatorname{diag}(1,1,1,-1,-1,-1)$, $Ge_i=e_{i+3}$, $Ge_{i+3}=-e_i$,
so $G^2=-\id$, $G^{\!\top}=-G$, $G^\dagger:=\eta^{-1}G^{\!\top}\eta=G$, $\eta(GX,GY)=-\eta(X,Y)$ (Proposition~3.3, Theorem~16.2 of \cite{ES17}), and
$G=-\cK\circ\Ad_{e_1e_2e_3}$, $\cK=\tfrac12\ad_{\mathsf P}$ (Lemma~3.4 of \cite{ES17}). For $Y\in V$ with components $(\bB;\bE)$ along
$(e_1,e_2,e_3;e_4,e_5,e_6)$, $z(Y)=\bB+i\bE$, so $G$ acts as $i$ (Theorem~12.3 of \cite{ES17}), $q(z)=z\cdot z=|\bB|^2-|\bE|^2+2i\,\bB\cdot\bE$, and
\[
E_6(Y):=\tfrac12|z(Y)|^2,\qquad E_0(Y):=\tfrac12|q(z(Y))|,\qquad \mathbf S(Y):=\tfrac i2\,\overline{z(Y)}\times z(Y)=\bE\times\bB
\]
are the undivided energy, the rest energy (the rest-frame energy density of Section~14.4 of \cite{ES17}) and the frame Poynting vector of $Y$; as
$\Btext=-\bB$ (Section~12 of \cite{ES17}), $\mathbf S=-\bE\times\Btext$ is minus the textbook one, which the rest-frame velocity of Section~14.4 uses;
only $|\mathbf S|$ enters below. A covector $p$ is a column, $C(p)=p^{\!\top}\eta^{-1}p$, $D(p)=p^{\!\top}G\eta^{-1}p$, $w(p):=C-iD=q(z(\pi))$, so
$E_0(\pi)=\tfrac12|w(p)|$, $E_6(\pi)=\tfrac12|p|^2$. Lorentz-class inertia (Theorem~17.11(ii) of \cite{ES17}) is written
$N=(\nu_1\eta^{-1}+\nu_2G\eta^{-1})/m$, $\nu:=\nu_1+i\nu_2$ (the $(a,b)$ of \cite{ES17}), as in \cref{d1:thm:closure}, keeping $b$ for the Yang--Mills
field. Units are $\hbar=c=1$ unless shown; $\tau$ is a time; $\mathrm K(\mu)$, $\mathrm E(\mu)$ are the complete elliptic integrals in the parameter $\mu$
\cite{DLMF}. d3-$n$ refers to \texttt{verify/d3\_closure\_dynamics.py} (integrations: eighth-order Runge--Kutta, relative tolerance $10^{-12}$; d3-0
checks the conventions), x3-$n$ to the independent \texttt{verify/x3\_check.py}.

\subsection{What the closure conserves}
\label{d3:sec:conserve}

The only $\lG$-equivariant closure in which the rotor rotates the frame covector without dilating it is $\mathfrak s=\R\hat g\ltimes_GV$,
$[\hat g,X]=GX$, $[X,Y]=0$ (Theorem~17.7 of \cite{ES17}). On $\mathfrak s^*$, with coordinates $(h,p)$, $F_h=\partial F/\partial h$ and
$\nabla F=\partial F/\partial p$, the Lie--Poisson equations of an arbitrary Hamiltonian $H$, and the bracket for which they read $\dot F=\{F,H\}$, are
\begin{equation}
\label{d3:eq:LPgen}
\dot p=H_h\,G^{\!\top}p,\qquad \dot h=-\,p^{\!\top}G\,\nabla H,\qquad
\{F_1,F_2\}=p^{\!\top}\bigl(F_{2,h}\,G\nabla F_1-F_{1,h}\,G\nabla F_2\bigr).
\end{equation}
For $H=h^2/(2I_0)+\tfrac12p^{\!\top}Np$ this is Proposition~17.8 of \cite{ES17}, $\dot p=(h/I_0)G^{\!\top}p$, $\dot h=-p^{\!\top}GNp$; re-derived from the
structure constants, with $H$ conserved for every symmetric $N$ and $|p|^2$ for every $H$ (d3-1).

\begin{theorem}[The momentum field is frozen up to its duality phase]
\label{d3:thm:frozen}
\begin{enumerate}[label=(\roman*),leftmargin=2.2em]
\item A function $f(p)$ is a Casimir of $\mathfrak s^*$ iff it is invariant under $p\mapsto e^{tG}p$, i.e.\ under the duality rotations
$z(\pi)\mapsto e^{-it}z(\pi)$ of the momentum field. So $E_6(\pi)=\tfrac12|p|^2$, $E_0(\pi)$ and the vector $\mathbf S(\pi)$ are Casimirs; $C$, $D$ are not.
\item At $p\neq0$ the symplectic leaves are two-dimensional, so there are five independent local Casimirs; the real and imaginary parts of the
$\bar z_jz_k$, $z=z(\pi)$, provide them (Jacobian rank five), and $E_6(\pi)$, $E_0(\pi)$, $\mathbf S(\pi)$ account for four, with
$E_6(\pi)^2=E_0(\pi)^2+|\mathbf S(\pi)|^2$.
\item Along the flow of any Hamiltonian, $p(\tau)=e^{-\theta(\tau)G}p_0$ with $\dot\theta=H_h$, $\theta(0)=0$, so
$z(\pi(\tau))=e^{i\theta(\tau)}z(\pi_0)$ and $w(p(\tau))=e^{2i\theta(\tau)}w(p_0)$.
\end{enumerate}
\end{theorem}

\begin{proof}
(i) For $f=f(p)$, \eqref{d3:eq:LPgen} gives $\{f,F\}=F_h\,p^{\!\top}G\nabla f=-F_h\,(Gp)\cdot\nabla f$, zero for every $F$ iff $f$ is constant along
$t\mapsto e^{tG}p$. As $\eta G=-G\eta$, $\eta^{-1}e^{tG}p=e^{-tG}\pi$, on which $G$ acts as $i$. $|z|^2$, $|z\cdot z|$ and $\bar z\times z$ are phase
invariant, while $C=\Re w$ and $D=-\Im w$ rotate into each other; $|\pi|=|p|$ as $\eta$ is orthogonal. (ii) $\{p_i,p_j\}=0$ and $\{h,p_j\}=(Gp)_j$, so
the Poisson tensor has rank two at $p\neq0$ and $\dim\mathfrak s^*-2=5$. The $\bar z_jz_k$ are phase invariant with Jacobian rank five at a generic
point; the relation is the Lagrange identity $u^2-|\mathbf S|^2=\tfrac14|q|^2$ of Section~14.3 of \cite{ES17}. (iii) $\dot p=-H_hGp$ integrates to
$p=e^{-\theta G}p_0$, so $\pi=e^{\theta G}\pi_0$, and $w$ is quadratic in $z$.
\end{proof}

The theorem is fixed, and in substance a corollary of Proposition~17.8 of \cite{ES17}, which identifies the coadjoint orbit through $(h,p_0)$ with
the cylinder $\R\times\{e^{\theta G}p_0\}$, so that every function of $p$ constant on these circles is a Casimir; (iii) is the phase law of
Theorem~17.13 of \cite{ES17} and the relation in (ii) is inherited. What is added is the naming and the count of the Casimirs. Already in the
clamped interior $p=e^{-\kappa\tau G}p_0$, so these functions of the momentum field were constant before the closure (\cref{d1:thm:closure}). All
items were checked exactly (d3-2); along five random Lorentz-class flows $H$, $|p|^2$, $E_6(\pi)$, $E_0(\pi)$, $\mathbf S(\pi)$ and (iii) hold to
$7\times10^{-12}$ (d3-3).

\begin{corollary}[What the closure can and cannot change]
\label{d3:cor:nogen}
\begin{enumerate}[label=(\roman*),leftmargin=2.2em]
\item For every Hamiltonian on $\mathfrak s^*$, quadratic or not, of Lorentz class or not, $E_0(\pi)$, $E_6(\pi)$ and $\mathbf S(\pi)$ of the momentum
field are constant along every solution: a null momentum field stays null, and one of rest energy $E_0$ keeps it.
\item The same functions of the frame field are not constants of motion. In the clamped split-form interior, $z(X)=c_-e^{-i\kappa\tau}+c_+e^{i\kappa\tau}$
(Proposition~16.5(ii) of \cite{ES17}) gives
\[
E_6(X)=\tfrac12\bigl(|c_-|^2+|c_+|^2\bigr)+\Re\bigl(\bar c_-\!\cdot c_+\,e^{2i\kappa\tau}\bigr),\quad
q(z(X))=c_-\!\cdot c_-\,e^{-2i\kappa\tau}+2\,c_-\!\cdot c_++c_+\!\cdot c_+\,e^{2i\kappa\tau},
\]
so $E_6(X)$ is constant iff $\bar c_-\!\cdot c_+=0$, and $c_-=(1,i,0)$, $c_+=(1,0,0)$ give $|q(z(X))|^2=9-8\sin^2\kappa\tau$. Under the free-rotor drift of
\cref{d3:ex:drift} with $x_0=0$, the frame field $X=e^{-\theta G}x$ has $E_0(X)=0$ at $\tau=0$ and $\tfrac12(3\,\Delta x)^2\approx1.807$ after three full
turns; on every rotating solution with $R_0>0$, $E_0(X)$ and $E_6(X)$ grow without bound.
\end{enumerate}
\end{corollary}

\begin{proof}
(i) Casimirs Poisson-commute with every Hamiltonian (\cref{d3:thm:frozen}(i)). (ii) Expand $|z|^2$ and $z\cdot z$; in the example $c_-\cdot c_-=0$,
$c_-\cdot c_+=c_+\cdot c_+=1$, so $|q|^2=|2+e^{2i\kappa\tau}|^2$. In \cref{d3:ex:drift} the motion is $2\pi$-periodic in $\theta$, so at $\theta=2\pi k$,
$e^{-\theta G}=\id$ and $x=k\,\Delta x$, a real multiple of $e_1$, whence $q(z(X))=k^2\Delta x^2$. The growth is \cref{d1:thm:closure}, item~4, from
Theorem~17.16 of \cite{ES17}.
\end{proof}

Item (i) is an obstruction for the momentum field, checked with the generic Hamiltonian
$h^2/(2I_0)+\tfrac1{10}h^4+\sin(p^{\!\top}M_1p)+\tfrac15h\,p^{\!\top}M_2p$, $M_{1,2}$ random symmetric: the Casimirs and $H$ stay constant to
$5\times10^{-12}$ while the phase of $w$ moves through $6.3$ radians (d3-4; x3-3). Item (ii) is fixed, from Proposition~16.5(ii) and Theorem~17.16 of
\cite{ES17} (d3-22; x3-4, x3-23). This first verdict cuts both ways. For the momentum field, $E_6(\pi)$ and its split
$E_6(\pi)^2=E_0(\pi)^2+|\mathbf S(\pi)|^2$ are labels of the leaf: its rest energy is carried by the dynamics, not produced by it. For the frame
field, to which the dictionary's $u=\tfrac12|z|^2$ and $m\leftrightarrow\tfrac12|q|$ refer, the closure does change the rest energy, by an
oscillation or an unbounded drift, not by a mechanism that selects a value; and $E_6(X)$ is not a function of $\tfrac12|p|^2=E_6(\pi)$ (on
zero-momentum orbits with $X\neq0$, $|p|^2=0<E_6(X)$, \cref{d2:prop:duality}). What moves is the duality angle $\theta$ and the rate $h/I_0$.

\subsection{The duality pendulum}
\label{d3:sec:pendulum}

From now on $N=(\nu_1\eta^{-1}+\nu_2G\eta^{-1})/m$ is of Lorentz class and $\nu\,w(p_0)=R_0e^{i\varphi_0}$, $R_0\ge0$. On a leaf $(\theta,h)$ are
canonical and $H=h^2/(2I_0)+U(\theta)$ (Proposition~17.8 of \cite{ES17}).

\begin{proposition}[The potential of the duality angle]
\label{d3:prop:pendulum}
$\tfrac12p^{\!\top}Np=\Re\bigl(\nu\,w(p)\bigr)/(2m)$, and with $\psi:=2\theta+\varphi_0$,
\[
U(\theta)=\frac{R_0}{2m}\cos\psi,\qquad R_0=|\nu|\,|w(p_0)|=2|\nu|\,E_0(\pi_0),\qquad I_0\ddot\theta=\dot h=\frac{R_0}m\sin\psi .
\]
\end{proposition}

\begin{proof}
$\tfrac12p^{\!\top}Np=(\nu_1C+\nu_2D)/(2m)=\Re(\nu(C-iD))/(2m)$, and $w(p(\theta))=e^{2i\theta}w(p_0)$ by \cref{d3:thm:frozen}(iii). For
$p(\theta)=e^{-\theta G}p_0$, $dU/d\theta=-p^{\!\top}NGp=p^{\!\top}GNp$ ($N$ symmetric, $G$ antisymmetric), so $\dot h=-dU/d\theta$; and $\dot\theta=h/I_0$.
\end{proof}

This restates Theorem~17.13 of \cite{ES17} (fixed there), re-verified exactly and along three random flows (d3-5); $R_0=2|\nu|E_0(\pi_0)$ only renames
$\tfrac12|w|$, and $E_6(\pi)$, $\mathbf S(\pi)$ do not enter.

\begin{proposition}[Small oscillations]
\label{d3:prop:omega}
The stable equilibrium is $\psi=\pi$, $U=-R_0/(2m)$. With $\psi=\pi+\xi$, $\ddot\xi=-\Omega^2\sin\xi$ and
\[
\Omega^2=\frac{2R_0}{mI_0}=\frac{4|\nu|\,E_0(\pi_0)}{mI_0};
\]
$\theta$ oscillates at the same frequency with half the amplitude, and at amplitude $0<\xi_0<\pi$ the period is $4\mathrm K(\sin^2(\xi_0/2))/\Omega$.
\end{proposition}

\begin{proof}
$\ddot\psi=2\ddot\theta=(2R_0/mI_0)\sin\psi$ and $\sin(\pi+\xi)=-\sin\xi$. The period follows from $\tfrac12\dot\xi^2=\Omega^2(\cos\xi-\cos\xi_0)$ with
$\sin(\xi/2)=\sin(\xi_0/2)\sin t$.
\end{proof}

Fixed, as a corollary of Theorem~17.13 of \cite{ES17}, whose pendulum it linearises; the period is the textbook one. At amplitude $10^{-3}$ the
measured period exceeds $2\pi/\Omega$ by $6.25\times10^{-8}$, the shift $\xi_0^2/16$; at amplitude $2$ it matches the elliptic formula to
$5\times10^{-14}$ (d3-6).

\begin{corollary}[When the rate is constant]
\label{d3:cor:flat}
\begin{enumerate}[label=(\roman*),leftmargin=2.2em]
\item If $\nu\neq0$, the rate $h/I_0$ is constant and nonzero iff $R_0=0$, that is $E_0(\pi_0)=0$ (a null momentum field, $|\bB|=|\bE|$, $\bB\cdot\bE=0$
\cite{Synge56}), and $h(0)\neq0$; then $U\equiv0$ and $\theta=h\tau/I_0$, an exact flat direction. For $E_0(\pi_0)>0$ the rate is constant only at
the equilibria $h(0)=0$, $\psi(0)\in\pi\Z$, where it vanishes and the frame moves at the constant velocity $Np_0$, $|Np_0|=|\nu||p_0|/m\neq0$.
\item If $\nu=0$, then $U\equiv0$ and the rate is constant for all data.
\item For duality-class inertia, $[N,G]=0$ (nine dimensions), the rate is constant for all data and $U=\tfrac12p_0^{\!\top}Np_0$ is constant, in
general nonzero.
\end{enumerate}
\end{corollary}

\begin{proof}
(i) $\dot h=(R_0/m)\sin\psi$. If $R_0>0$ and $h$ is constant, $\sin\psi\equiv0$, so $\theta$ is constant and $h=I_0\dot\theta=0$; conversely these points
are equilibria, where Theorem~17.12 of \cite{ES17} ($N_c=0$, $\theta(0)=0$) gives $\dot x=Np_0$, and $N^2=|\nu|^2/m^2$. (ii) $N=0$. (iii) $GN$ is
antisymmetric, so $\dot h=-p^{\!\top}GNp=0$, and $U(\theta)=\tfrac12p_0^{\!\top}e^{\theta G}Ne^{-\theta G}p_0=\tfrac12p_0^{\!\top}Np_0$.
\end{proof}

Items (i)--(ii) restate Corollary~17.14 of \cite{ES17} (fixed there), with the equilibria explicit as in \cref{d1:thm:closure}; (iii) restates
Theorems~17.11(i),(iii) and~17.12 of \cite{ES17}. Numerically $|h-h_0|<10^{-14}$ over $\tau\in[0,30]$ on null data and for a random duality-class $N$,
and at the equilibria $h$ stays $0$ while $x=\tau Np_0$ (d3-7).

\begin{remark}[A pseudo-Goldstone reading, and its limits]
\label{d3:rem:goldstone}
On a leaf $\theta$ is the angle of the duality group $U(1)_G$ of the momentum field, broken by a nonzero rest energy (\cref{d2:prop:duality}): in
the duality class it is a symmetry of $H$, in the Lorentz class it is broken on each leaf in proportion to $E_0(\pi)$. Exactly (d3-8),
\[
I_0\,\Omega^2=U''(\theta_{\min})=\frac{2R_0}m=\frac{4|\nu|\,E_0(\pi_0)}m .
\]
The first equality defines $\Omega$ for every one-degree-of-freedom oscillator; the closure-specific content is the second, $U''=2R_0/m$, which is
Theorem~17.13 of \cite{ES17}. With $I_0$ for $f_\pi^2$ and $|\nu|E_0(\pi_0)/m$ for the explicit breaking it has the shape of the
Gell-Mann--Oakes--Renner relation \cite{GellMannOakesRenner68,Weinberg72}, as every pendulum does; the reading is interpretive. The breaking is
explicit, by the inertia, not spontaneous: there is no vacuum manifold, $E_0(\pi)$ is a label fixed by the data (\cref{d3:cor:nogen}), not an order
parameter (Section~14.3 of \cite{ES17}), and with one degree of freedom ``mass'' means a frequency. What survives: on a null momentum field the angle
is an exact zero mode, and a rest energy of the momentum field gives it the frequency $\Omega\propto\sqrt{E_0(\pi_0)}$.
\end{remark}

\subsection{The energy budget: a rate part and a frame part}
\label{d3:sec:budget}

\begin{proposition}[Rate part and frame part]
\label{d3:prop:budget}
With Lorentz-class inertia, along every solution:
\begin{enumerate}[label=(\roman*),leftmargin=2.2em]
\item $H=K_{\rm rot}+U$, $K_{\rm rot}=\tfrac12I_0\kappa^2\ge0$, $\kappa=h/I_0$, $U=\tfrac{R_0}{2m}\cos\psi=\Re(\nu\,w(p))/(2m)$, and
$\dot K_{\rm rot}=-\dot U=\tfrac{R_0}m\,\kappa\sin\psi$;
\item in the rotating regime $K_{\rm rot}$ runs through exactly $[H-\tfrac{R_0}{2m},H+\tfrac{R_0}{2m}]$ in every period of the rate, so the energy exchanged
is $R_0/m=2|\nu|E_0(\pi_0)/m$; in the librating regime it runs through $[0,H+\tfrac{R_0}{2m}]$;
\item $K_{\rm rot}$ and $U$ are separately invariant under $\lG$ acting by $(h,p)\mapsto(h,\Lambda^{-\top}p)$, $\Lambda=e^r$, $r\in\lG$, and the flow
commutes with it;
\item $(\theta,h)(\tau)$ depends on $p_0$ only through $\nu\,w(p_0)$: $E_6(\pi_0)$ and $\mathbf S(\pi_0)$ do not enter the motion.
\end{enumerate}
\end{proposition}

\begin{proof}
(i) By \cref{d3:prop:pendulum}, $\dot K_{\rm rot}=h\dot h/I_0$ and $\dot U=-(R_0/2m)\,2\kappa\sin\psi$. (ii) Rotating, $\dot\psi=2\kappa$ never vanishes and
$\cos\psi$ covers $[-1,1]$ in each period; librating, $K_{\rm rot}=0$ at the turning points and is largest at $\psi=\pi$. (iii) The Lorentz requirement
$rN+Nr^{\!\top}=0$ (Definition~17.10 of \cite{ES17}) integrates to $\Lambda^{-1}N\Lambda^{-\top}=N$, so $U(\Lambda^{-\top}p)=U(p)$; $\Lambda^{-\top}$ commutes with
$G^{\!\top}$ and maps solutions to solutions. (iv) $\pi\mapsto\Lambda\pi$ as $\Lambda\in O(3,3)$, and $\lG\cong\so{3,\C}$ preserves $q$ (Theorem~12.1 of
\cite{ES17}), so $w(\Lambda^{-\top}p)=w(p)$; the equations for $(\theta,h)$ contain only $R_0$ and $\varphi_0$.
\end{proof}

Fixed; (i) restates the energy $E_\theta=H$ of Theorem~17.13 of \cite{ES17}, (ii)--(iv) are added. Numerically $H=K_{\rm rot}+U$ to $3\times10^{-11}$,
the range of $K_{\rm rot}$ equals $R_0/m$ to $3\times10^{-10}$, and a Lorentz transformation of $p_0$ changed $E_6(\pi_0)$ by $5.0$ and $\theta(\tau)$ by
less than $10^{-11}$ (d3-9). So the conserved total differentiates into two Lorentz-invariant parts whose exchange is fixed by the rest energy of the
momentum field alone; neither is the undivided $E_6(\pi)$, which is inert.

\paragraph{Units, and the factor one half.} With $[\tau]=\mathsf T$, $[\kappa]=\mathsf T^{-1}$ (Section~3 of \cite{ES17}) and $[X]=\mathsf L$, a
convention by \cref{d3:prop:scalings}(ii) that makes $m$ a mass (d3-10): $[I_0]=\mathsf{ML}^2$, $[h]=\mathsf{ML}^2\mathsf T^{-1}$ (an action),
$[p]=\mathsf{MLT}^{-1}$, $[N]=\mathsf M^{-1}$, $[H]=[K_{\rm rot}]=[U]=\mathsf{ML}^2\mathsf T^{-2}$, $[M_{\rm rot}]=\mathsf M$ for $M_{\rm rot}:=I_0/\ell^2$; and
$E_0(\pi)$ is a momentum squared, $|\nu|E_0(\pi)/m$ an energy. With the clamped value $\kappa=c/\ell$,
$K_{\rm rot}=\tfrac12I_0\kappa^2=\tfrac12M_{\rm rot}c^2$, the ``$c^2$ part'' of the conserved total (interpretive): literally the Newtonian kinetic energy
of a mass $M_{\rm rot}$ on a circle of radius $\ell$ at rim speed $c$, outside Newtonian kinetics. It is not $E=mc^2$: the $\tfrac12$ comes from an
energy quadratic in the rate, where Planck's relation is linear; $M_{\rm rot}$ is not tied to $m$ (\cref{d3:prop:scalings}(iv)); and $\kappa=c/\ell$
holds only for the clamped rotor.

\subsection{When the rate is constant}
\label{d3:sec:multiscale}

\begin{proposition}[The multiscale criterion]
\label{d3:prop:multiscale}
Let $R_0>0$, $\kappa_0=\kappa(0)\neq0$, $\psi_0=\psi(0)$ and $\varepsilon:=R_0/(mI_0\kappa_0^2)=U_{\max}/K_{\rm rot}(0)$, $U_{\max}=R_0/(2m)$. Then
$\kappa^2=\kappa_0^2[1+\varepsilon(\cos\psi_0-\cos\psi)]$, so $|\kappa^2/\kappa_0^2-1|\le2\varepsilon$; the rotor turns for all time iff
$\varepsilon(1-\cos\psi_0)<1$, which holds for every phase when $\varepsilon<\tfrac12$; for $\varepsilon(1-\cos\psi_0)>1$ it librates, and equality is
the separatrix.
\end{proposition}

\begin{proof}
Divide the conservation of $H=\tfrac12I_0\kappa^2+\tfrac{R_0}{2m}\cos\psi$ by $\tfrac12I_0\kappa_0^2$. The least value of $\kappa^2$, at $\cos\psi=1$, is
$\kappa_0^2[1-\varepsilon(1-\cos\psi_0)]$; if it is negative, $\kappa$ vanishes where $\sin\psi\neq0$, hence $\ddot\theta\neq0$, and changes sign.
\end{proof}

This restates Corollary~17.15 of \cite{ES17} (fixed there), checked along four flows to $8\times10^{-12}$ and in five cases on both sides of the
criterion (d3-11); the form $\varepsilon=U_{\max}/K_{\rm rot}(0)$ is ours. The reading is interpretive (d3-12): if $c$ is the fixed conversion of
Theorem~3.2 of \cite{ES17}, a varying $\kappa=c/\ell$ is a varying length $\ell(\tau)=c/\kappa(\tau)$, with relative spread
$\varepsilon+O(\varepsilon^2)$ along a rotating solution (measured at $\varepsilon=10^{-1},10^{-2},10^{-3}$). A constant
``$c/\ell$'' requires $\tfrac12M_{\rm rot}c^2\gg|\nu|E_0(\pi_0)/m$: the rotor energy must dominate the rest-energy amplitude of the momentum field. The
clamped split-form interior of Section~16 of \cite{ES17} is the limit $\varepsilon\to0$ (Proposition~17.9 and Corollary~17.15 of \cite{ES17}); the
positive-form interior of Section~15 has duality-class inertia and an exactly constant rate (Theorem~17.11(iii) of \cite{ES17}). The ``symmetric''
$c^2/\ell^2$ is constant only insofar as the differentiated part is negligible beside the rotor, and $\varepsilon$ is an input.

\subsection{The rest energy as a rate of rotation}
\label{d3:sec:rate}

\begin{proposition}[Two rotations at rest]
\label{d3:prop:tworot}
(i) For the real first-order Dirac form of Section~4.4 of \cite{ES17} at $\bp=0$, $H(0)=me_1$, $Je_1=T=e_1e_2e_3$, $T^2=-\id$, and
$\Psi(\lambda)=\cos(m\lambda)\Psi_0-\sin(m\lambda)Je_1\Psi_0=e^{-m\lambda T}\Psi_0$: angular frequency $m$, that is $mc^2/\hbar$. (ii) The zero-momentum
orbits $X(\tau)=e^{-\kappa G\tau}X_0$ satisfy $\alpha_+X=0$, $\ddot X=-\kappa^2X$: angular frequency $\kappa=c/\ell$, $z\mapsto e^{-i\kappa\tau}z$,
$q\mapsto e^{-2i\kappa\tau}q$. The orbits $X(\tau)=e^{\kappa G\tau}X_0$, $z=e^{i\kappa\tau}c_+$, are duality rotations of the other sense at the same rate,
with momentum $\alpha_+X=2\kappa GX\neq0$.
\end{proposition}

\begin{proof}
(i) $\dot\Psi=-JH\Psi$ with $JH(0)=mT$, since $e_1$ commutes with $J=e_2e_3$. (ii) $\tfrac{d}{d\tau}e^{\mp\kappa G\tau}=\mp\kappa Ge^{\mp\kappa G\tau}$,
$G^2=-\id$.
\end{proof}

Item (i) is inherited, the realification of $\psi=e^{-imc^2t/\hbar}\psi_0$, de Broglie's internal clock \cite{deBroglie1925} (compare
\cref{d2:prop:restflow}); (ii) restates Proposition~16.5(i)--(ii) of \cite{ES17} (fixed there); both checked exactly (d3-13). Duality acts along both
senses at the same rate: the zero-momentum sense is singled out by the momentum, not by duality.

\begin{proposition}[No operator relates the two rates]
\label{d3:prop:norel}
(i) $\Ad_{e^{\vartheta T}}$ fixes $e_1,e_2,e_3$ and maps a spatial $e_a$ to $\cos2\vartheta\,e_a+\sin2\vartheta\,Te_a$, a vector only for
$\vartheta\in\tfrac\pi2\Z$: the Dirac rest flow is not a motion of $V$, and $T$, of grade three, is not in $\so{3,3}$. (ii)
$(e^{\theta G})^{\!\top}\eta\,e^{\theta G}=\cos2\theta\,\eta+\sin2\theta\,\eta G$, so $e^{\theta G}\in O(3,3)$, and lifts to $\mathrm{Pin}(3,3)$, only for
$\theta\in\pi\Z$. Hence neither one-parameter group is the image of the other under the spin homomorphism, so this route supplies no relation
between $m$ and $\kappa$, which remain independent inputs.
\end{proposition}

\begin{proof}
(i) $T$ commutes with the temporal and anticommutes with the spatial generators (Lemma~2.1 of \cite{ES17}), so
$e^{\vartheta T}e_ae^{-\vartheta T}=e^{2\vartheta T}e_a$, and $Te_a$ has grade four. (ii) Expand with $e^{\theta G}=\cos\theta+\sin\theta\,G$,
$G^{\!\top}\eta=\eta G$ symmetric, $G^{\!\top}\eta G=-\eta$.
\end{proof}

Item (ii) is the remark of Sections~14.3 and~17.6 of \cite{ES17}. The proposition obstructs one route only, a map of the arena from one rate to the
other (d3-15); $m$ and $\kappa$ are independent because both are free inputs. What remains is an added identification, interpretive (d3-14): with
$E=\hbar\times$rate, $\hbar\kappa=mc^2$ holds iff $\ell=\hbar/(mc)$, the reduced Compton wavelength. The dynamics does not supply it. In the interior
$\ddot X=-\kappa^2X$ does not contain $m$; and the rescaling $(\tau,\kappa,I_0,m)\mapsto(s\tau,\kappa/s,sI_0,sm)$ of \cref{d3:prop:scalings}, which maps
solutions to solutions and preserves the action (hence any quantisation of $h$), keeps $m\kappa$ and multiplies $m/\kappa$ by $s^2$, so
$m=\hbar\kappa/c^2$ picks one point of that orbit by fiat. The rate is also ambiguous by a factor two ($q$ turns at $2\kappa$, $z$ at $\kappa$, the Dirac
bilinears at $0$ or $2m$): $2\hbar\kappa=mc^2$ gives $\ell=2\hbar/(mc)$, and the algebra decides neither.

\subsection{Scale freedom}
\label{d3:sec:scale}

\begin{proposition}[Scalings]
\label{d3:prop:scalings}
\begin{enumerate}[label=(\roman*),leftmargin=2.2em]
\item $\eta$, $G$, $G\eta^{-1}$, $\eta G$, $\cK$ and the structure constants of the complex cross product $\beta_0$ are integer matrices.
\item $\ddot X=-\kappa^2X$ is invariant under $X\mapsto\sigma X$, which scales $\tfrac m2\eta(\alpha_+X,\alpha_+X)$ by $\sigma^2$, absorbed by
$m\mapsto m/\sigma^2$, and under $(\tau,\kappa)\mapsto(s\tau,\kappa/s)$, which with $m\mapsto sm$ preserves the action.
\item The closure is invariant under $(\tau,I_0,m)\mapsto(s\tau,sI_0,sm)$ with $(\theta,h,p)$ unchanged as functions of the rescaled time; $\kappa\mapsto\kappa/s$,
and $m\kappa$, $\varepsilon$ and the action are invariant.
\item For $s,t>0$, $(I_0,m)\mapsto(sI_0,tm)$ at fixed initial rate $\kappa_0$ ($h_0\mapsto sh_0$) keeps $I_0>0$ and $N$ of Lorentz class and multiplies
$\varepsilon$ by $1/(st)$ (by $s/t$ at fixed $h_0$).
\end{enumerate}
\end{proposition}

\begin{proof}
(i) $\beta_0(e_i,e_j)=\epsilon_{ijk}e_k$, $\beta_0(e_i,e_{j+3})=\epsilon_{ijk}e_{k+3}$, $\beta_0(e_{i+3},e_{j+3})=-\epsilon_{ijk}e_k$; the rest is Lemma~3.4 of
\cite{ES17}. (ii) Linearity; $Y(\tau')=X(\tau'/s)$ solves $Y''=-(\kappa/s)^2Y$, the Lagrangian scales by $s^{-1}$ and $d\tau'=s\,d\tau$. (iii) With
$d/d\tau'=s^{-1}d/d\tau$, $I_0'=sI_0$, $N'=N/s$, both equations of Proposition~17.8 of \cite{ES17} are reproduced, $H'=H/s$, and $\kappa_0'=\kappa_0/s$ gives
$\varepsilon'=\varepsilon$. (iv) $R_0$ does not depend on $(I_0,m)$, and $\varepsilon=R_0/(mI_0\kappa_0^2)$ with $\kappa_0$ held fixed; at fixed $h_0$,
$\kappa_0\mapsto\kappa_0/s$.
\end{proof}

Items (i) and (iv) restate the remark closing Section~17.5 of \cite{ES17} (fixed there; the fixed-$\kappa_0$ qualifier is item~12 of
\texttt{investigations/errata\_2026-10-17}); (ii) and (iii) are added; all are fixed. The closure rescaled by $s=2.5$ reproduces $(\theta,h,p)$ to
$7\times10^{-12}$, and $\varepsilon$ measured from the range of $\kappa^2$ goes from $0.27209155$ to $0.04534858=\varepsilon/6$ under
$(I_0,m)\mapsto(2I_0,3m)$ at fixed $\kappa_0$, and to $\tfrac23\varepsilon$ at fixed $h_0$ (d3-16, d3-17).

\begin{proposition}[Nothing in the algebra fixes $\ell$]
\label{d3:thm:noscale}
(i) Schur's lemma \cite{Schur} fixes only the form $\mathcal C=c\,\id$ of the conversion operator (Theorem~3.2 of \cite{ES17}), and $c$ converts units of time into
units of length. (ii) No monomial $c^{n_1}\hbar^{n_2}$ is a length, a mass or a time; with a mass $m$ adjoined, the only length is $\hbar/(mc)$ up to a
number. (iii) The interior systems with different $\kappa>0$ are carried into each other by a rescaling of $\tau$ that fixes every Clifford datum
(\cref{d3:prop:scalings}(ii)); so no condition formulated in the Clifford data distinguishes one value of $\kappa$, or of $\ell=c/\kappa$, and in the
closure none fixes $\varepsilon$ (\cref{d3:prop:scalings}(iv)).
\end{proposition}

\begin{proof}
(i) The commutant of the vector representation of $\so3$ in $M_3(\R)$ is $\R\,\id$. (ii) In the basis $(\mathsf M,\mathsf L,\mathsf T)$,
$[c]=(0,1,-1)$ and $[\hbar]=(1,2,-1)$ span a plane containing none of $(0,1,0)$, $(1,0,0)$, $(0,0,1)$; adjoining $[m]=(1,0,0)$ gives an invertible
matrix, with $(0,1,0)=-[c]+[\hbar]-[m]$. (iii) The data are integer matrices, and the rescaling acts on $\tau$ and $\kappa$ only.
\end{proof}

This is inherited dimensional analysis, already implicit in equation~(kappa) of Section~3 of \cite{ES17} (``their sum requires a scale'') and in the
remark closing Section~17.5 (d3-18). It is the one exact form of ``$m$ breaks the symmetry of $c^2$'' available here: $c$ carries no scale, the shift
$\kappa G$ needs a length $\ell$ that the algebra cannot supply, and through $\ell=\hbar/(mc)$ that length is a mass. What it ``breaks'' is the change
of the unit of $\tau$ (\cref{d3:prop:scalings}(iii)), not a symmetry of a vacuum; the reading is interpretive.

\subsection{The frame drifts}
\label{d3:sec:drift}

\begin{example}[Example~17.17 of \cite{ES17}, re-verified]
\label{d3:ex:drift}
Take $p_0=e^1$, $m=I_0=\kappa_0=1$, $\nu=1$: $w=1$, $\varphi_0=0$, $\varepsilon=1$, $\dot\theta^2=1+2\sin^2\theta$, and by Theorem~17.12 of \cite{ES17}
($N_c=0$) $\dot x=e^{2\theta G}Np_0=\cos2\theta\,e_1+\sin2\theta\,e_4$. The $e_4$ part integrates to zero ($\theta\mapsto\pi-\theta$), and with
$(1-2\sin^2\theta)/\sqrt{1+2\sin^2\theta}=2/\sqrt{1+2\sin^2\theta}-\sqrt{1+2\sin^2\theta}$,
\[
\Delta x\big|_{\theta:0\to\pi}=\int_0^\pi\frac{\cos2\theta\,d\theta}{\sqrt{1+2\sin^2\theta}}\,e_1=\bigl(4\mathrm K(-2)-2\mathrm E(-2)\bigr)e_1\approx0.3168040511\,e_1 .
\]
Over a full turn, $\theta:0\to2\pi$, $\Delta x=(8\mathrm K(-2)-4\mathrm E(-2))e_1\approx0.6336081022\,e_1$, the value of Example~17.17 of \cite{ES17}, whose
corrected wording (erratum item~13) makes the turn $\theta\mapsto\theta+2\pi$, two periods of the rate. After half a turn $h$ has returned and
$p=-p_0$ ($e^{\pi G}=-\id$); after a full turn both have returned.
\end{example}

Fixed, as a re-verification (d3-19): quadratures agree with the elliptic expressions to $10^{-25}$, and integrating $(\theta,h,p,x)$ to the events
$\theta=\pi,2\pi$ reproduces both values and the returns of $h$, $p$ to $3\times10^{-14}$. By Theorem~17.16 of \cite{ES17} every rotating solution with a
nonzero rest energy of the momentum field has an unbounded frame; with a free rotor, a constant nonzero rate and a bounded frame coexist exactly on null
data (end of Section~17.5 of \cite{ES17}; at the equilibria of \cref{d3:cor:flat} the rate is zero and the frame moves uniformly), or, for several
constituents, on the balance surface (Proposition~17.18 of \cite{ES17}). In motion the differentiation unbinds the frame, and the frame field's rest
energy grows with it (\cref{d3:cor:nogen}(ii)).

\subsection{A quantum rotor, conditionally}
\label{d3:sec:quantum}

\begin{proposition}[The half turn]
\label{d3:prop:halfturn}
(i) $e^{\pi G}=-\id$, $e^{2\pi G}=\id$, and for $p_0\neq0$ the leaf circle $\{e^{\theta G}p_0\}$ has period exactly $2\pi$. (ii) For every symmetric $N$,
$U(\theta+\pi)=U(\theta)$. (iii) $\Gamma=e_1e_2\cdots e_6$ anticommutes with every $e_a$ and $\Gamma^2=\id$, so $\Ad_\Gamma=-\id$ on $V$: $\Gamma$ implements
$e^{\pi G}=-\id$ on spinors; by \cref{d3:prop:norel}(ii) the multiples of the half turn are the only rotor elements that lift to $\mathrm{Pin}(3,3)$.
\end{proposition}

\begin{proof}
(i) $(e^{\theta G}-\id)p_0=(\cos\theta-1)p_0+\sin\theta\,Gp_0$ with $p_0\perp Gp_0$, $|Gp_0|=|p_0|$. (ii) $p(\theta+\pi)=-p(\theta)$ and $U$ is quadratic.
(iii) A product of all six anticommuting generators anticommutes with each, and $(e_1\cdots e_6)^2=(-1)^{15}\prod_a\eta_{aa}=\id$.
\end{proof}

This restates the remark of Section~17.6 of \cite{ES17} (fixed there; in the realisation $\Gamma=J\,e_1e_4e_5e_6=+e_1\cdots e_6$), re-verified exactly
(d3-20). The quantum statement it permits is conditional on an assumption.

\begin{assumption}[A quantised leaf]
\label{d3:ass:quant}
The leaf $T^*S^1$ of \cref{d3:thm:frozen}(ii) is quantised by wave functions with $\Phi(\theta+2\pi)=e^{2\pi i\delta}\Phi(\theta)$, $h=-i\hbar\,\partial_\theta$,
and the Casimirs of the leaf are treated as numbers.
\end{assumption}

\begin{proposition}[A conditional rate ladder]
\label{d3:prop:ladder}
Under \cref{d3:ass:quant}: (i) $h$ has spectrum $\hbar(\Z+\delta)$; on flat leaves (\cref{d3:cor:flat}) the levels are
$E^{\rm rot}_n=\hbar^2(n+\delta)^2/(2I_0)$ (plus the constant $\tfrac12p_0^{\!\top}Np_0$ on duality-class leaves), $\kappa_n=(n+\delta)\hbar/I_0$ and
$E^{\rm rot}_{n+1}-E^{\rm rot}_n=\tfrac\hbar2(\kappa_n+\kappa_{n+1})$. (ii) If $R_0>0$ the Hamiltonian is a Mathieu operator \cite{Mathieu,DLMF} whose
potential couples $e^{i(n+\delta)\theta}$ only to $e^{i(n\pm2+\delta)\theta}$, so for every $\delta$ the parity of $n$ is conserved; it labels the
eigenvalues $\pm e^{i\pi\delta}$ of the half-turn translation $\theta\mapsto\theta+\pi$ ($\pm1$ for $\delta=0$). For $q_{\rm M}:=I_0R_0/(2m\hbar^2)\gg1$ the
lowest gap in either parity sector is $\hbar\Omega-\hbar^2/(2I_0)+O(q_{\rm M}^{-1/2}\hbar^2/I_0)$, and the two sectors are degenerate up to tunnelling.
(iii) If moreover $\delta=0$, and, as a separate assumption, the first level $h=\hbar$ is matched by $\hbar\kappa_1=mc^2$ with $\ell=\hbar/(mc)$, then
$I_0=\hbar^2/(mc^2)$, $M_{\rm rot}=m$, and $E^{\rm rot}_1=\tfrac12mc^2$.
\end{proposition}

\begin{proof}
(i) Spectrum of $-i\hbar\partial_\theta$ on twisted periodic functions; for $[N,G]=0$, $U$ is the constant of \cref{d3:cor:flat}(iii). (ii)
$\cos(2\theta+\varphi_0)$ shifts $n$ by $\pm2$, and $e^{i(n+\delta)(\theta+\pi)}=(-1)^ne^{i\pi\delta}e^{i(n+\delta)\theta}$. Dividing the Schr\"odinger equation by
$-\hbar^2/(2I_0)$ gives $\Phi''+(a_{\rm M}-2q_{\rm M}\cos(2\theta+\varphi_0))\Phi=0$, $a_{\rm M}=2I_0E/\hbar^2$, with
$a_{\rm M}\sim-2q_{\rm M}+2(2r_{\rm M}+1)\sqrt{q_{\rm M}}-\tfrac18((2r_{\rm M}+1)^2+1)$ as $q_{\rm M}\to\infty$ (DLMF~28.8.1); the levels $r_{\rm M}=0,1$ lie in the same sector and
differ by $4\sqrt{q_{\rm M}}-1$, and $\hbar\Omega=(\hbar^2/2I_0)\,4\sqrt{q_{\rm M}}$. (iii) $\hbar/I_0=c/\ell=mc^2/\hbar$ and $E^{\rm rot}_1=\hbar^2/(2I_0)$.
\end{proof}

Checked (d3-21; independently, x3-21): at $R_0=0$ the ladder is exact for $\delta=0$ and $0.3$; at $R_0/m=2000$, $\hbar=I_0=1$, the lowest gap in each
sector is $62.7415$ for both $\delta$, against $\hbar\Omega-\hbar^2/(2I_0)=62.7456$ ($62.7416$ with the next term of DLMF~28.8.1), and the sectors' ground
states agree to $3\times10^{-12}$. The proposition is constructed: the quantisation of the leaf exists (the orbit method, $\delta$ a character of the
stabiliser $2\pi\Z$ in the universal cover of the rotor group), and nothing selects $\delta$, $I_0$ or the matchings; whether it is physical is open.
Missing are a reason for the assumption (a quantum theory of the momentum field, which the closure does not give and the split-form frame obstructs,
having no ground state, \cref{d2:cor:noground}), a value of $I_0$, a coupling to matter (only the $\Z_2$ of the half turn reaches spinors, as $\Gamma$,
and nothing ties the parity of $n$ to chirality), and any derivation of the two matchings of (iii). $E^{\rm rot}_1=\tfrac12mc^2$ is
$K_{\rm rot}=\tfrac12h\kappa$ at $h=\hbar$: the $\tfrac12$ is built in, and $mc^2$ appears only as $\hbar$ times the rate put in, neither as a level nor
as a transition energy.

\subsection{Assessment}
\label{d3:sec:assess}

\begin{table}[t]
\centering
\small
\begin{tabular}{@{}>{\raggedright\arraybackslash}p{2.5cm}>{\raggedright\arraybackslash}p{8.3cm}>{\raggedright\arraybackslash}p{2.8cm}>{\raggedright\arraybackslash}p{1.1cm}@{}}
\toprule
claim & statement & class & check\\
\midrule
\cref{d3:thm:frozen} & $E_6(\pi)$, $E_0(\pi)$, $\mathbf S(\pi)$ are Casimirs; only the duality phase moves & fixed (cf.\ Prop.~17.8) & d3-1--3\\
\cref{d3:cor:nogen} & (i) no Hamiltonian changes $E_0(\pi)$; (ii) $E_0(X)$ oscillates or grows & (i) obstructed, (ii) fixed & d3-4, 22\\
\cref{d3:prop:pendulum} & $U=\tfrac{R_0}{2m}\cos(2\theta+\varphi_0)$, $R_0=2|\nu|E_0(\pi_0)$ & fixed (restated) & d3-5\\
\cref{d3:prop:omega} & $\Omega^2=4|\nu|E_0(\pi_0)/(mI_0)$; elliptic period & fixed (corollary) & d3-6\\
\cref{d3:cor:flat} & constant rate: null data, $\nu=0$, equilibria, duality class & fixed (restated) & d3-7\\
\cref{d3:rem:goldstone} & $U''=2R_0/m$ / $\theta$ as a pseudo-Goldstone angle & restated / interpretive & d3-8\\
\cref{d3:prop:budget} & $H=K_{\rm rot}+U$, exchange $R_0/m$, Lorentz invariance, $E_6(\pi)$ inert & fixed & d3-9\\
units & $K_{\rm rot}=\tfrac12M_{\rm rot}c^2$ is the ``$c^2$ part''; not $E=mc^2$ & interpretive & d3-10\\
\cref{d3:prop:multiscale} & $\kappa^2$ law, bound $2\varepsilon$, rotation criterion & fixed (restated) & d3-11\\
heavy rotor & constant $c/\ell$ needs $\tfrac12M_{\rm rot}c^2\gg|\nu|E_0(\pi_0)/m$ & interpretive & d3-12\\
\cref{d3:prop:tworot} & rest energy as a rotation rate (Dirac; frame, both senses) & inherited / restated & d3-13\\
matching & $\hbar\kappa=mc^2\iff\ell=\hbar/(mc)$, an added identification & interpretive & d3-14\\
\cref{d3:prop:norel} & no Spin or $O(3,3)$ map relates the two rates & obstructed (one route) & d3-15\\
\cref{d3:prop:scalings} & integer data; scalings of $X$, $\tau$, $I_0$, $m$ & fixed & d3-16, 17\\
\cref{d3:thm:noscale} & nothing in the algebra fixes $\ell$ (a mass scale) & inherited & d3-18\\
\cref{d3:ex:drift} & $\Delta x=(8\mathrm K(-2)-4\mathrm E(-2))e_1$ per full turn & fixed (restated) & d3-19\\
\cref{d3:prop:halfturn} & $e^{\pi G}=-\id$, implemented by $\Gamma$ & fixed (restated) & d3-20\\
\cref{d3:prop:ladder} & conditional ladder $\kappa_n=(n+\delta)\hbar/I_0$; $E^{\rm rot}_1=\tfrac12mc^2$ & constructed & d3-21\\
$b$-field & compact transports (\cref{d1:sec:main,d4:sec:bfield}): Zeeman, Jordan, screening; $\so4_\Delta$ & fixed (restated) & d3-23\\
\bottomrule
\end{tabular}
\caption{The claims of \cref{d3:sec:main}, their classes and the checks of \texttt{d3\_closure\_dynamics.py} (d3-0 verifies the conventions).
``Fixed (restated)'' and ``restated'': fixed in \cite{ES17} (or in another section of this note), restated and re-verified here.}
\label{d3:tab:claims}
\end{table}

Against the working dictionary (\cref{d3:tab:claims}) the closure gives four answers. \emph{$E$ as the Casimir $|p|^2$}: $E_6(\pi)=\tfrac12|p|^2$ is a
definition; the identification of $u=\tfrac12|z|^2$ of the frame with $\tfrac12|p|^2$ fails, since the frame drifts and its $u$ is not conserved
(\cref{d3:cor:nogen}(ii)); and the Casimir does not enter the rotor's motion. What is conserved and differentiates in motion is $H=K_{\rm rot}+U$.
\emph{$c^2$ as the factorising shift $\kappa G$}: in the closure $c^2$ appears as the rotor term $\tfrac12(I_0/\ell^2)c^2$ and, clamped, as
$(\kappa G)^2=-c^2/\ell^2$; $c$ is fixed up to units by Schur's lemma, $\ell$ is free, $\kappa$ is dynamical and constant only when $\varepsilon\ll1$, and
the factor $\tfrac12$ is that of a rotor. \emph{$m$ as the rest energy $\tfrac12|q|$}: confirmed in a precise dynamical sense for the momentum
field, whose $E_0(\pi)$ is the amplitude of the duality potential, sets $\Omega^2$ and the energy exchanged, and vanishes on the flat direction; but
there it is a Casimir, carried and not generated. For the frame field $\tfrac12|q(X)|$ is not conserved even in the clamped interior
(Proposition~16.5(ii) of \cite{ES17}) and grows without bound under the free-rotor drift, which changes it without selecting a value; and
$mc^2=\hbar\kappa$, which would tie $m$ to the rotor, is an added identification (\cref{d3:prop:norel}). \emph{$b$, the Yang--Mills field}: the closure
contains none and needs none (its frame is abelian, Theorems~17.6--17.7 of \cite{ES17}); a Yang--Mills field can only be added as a
transport. An admissible one is compact (Theorem~2.3 and Corollary~2.4 of \cite{ES17}): $b=(A_{\rm d}+A_{\rm a},A_{\rm d}-A_{\rm a})$ on the
temporal and spatial triples, rotation vectors $\mathbf a_{\rm d},\mathbf a_{\rm a}$, diagonal part in $\lG$, anti-diagonal part in $\mathfrak p_G$. In the
interior with the shift $b+\kappa G$ (\cref{d4:sec:bfield}) the diagonal part alone gives the bounded lines $\kappa$, $\kappa\pm|\mathbf a_{\rm d}|$, removable along a
worldline \cite{Larmor1897} and so not a gauge-invariant spectrum; the anti-diagonal part alone gives Jordan blocks at the rate
$\sqrt{\kappa^2+|\mathbf a_{\rm a}|^2}$ and secular growth; $\mathbf a_{\rm d}\cdot\mathbf a_{\rm a}=0$, $\mathbf a_{\rm d}\neq0$ screens it, every orbit bounded
although $[b,G]\neq0$, except on the ratio $|\mathbf a_{\rm d}|^2=\kappa^2+|\mathbf a_{\rm a}|^2$, where a static mode produces a Jordan block
(\cref{d4:thm:screen}); a generic compact $b$ is unbounded (d3-23). At three constituents $\so4_\Delta$ is the part of $\mathfrak{ps}_W$ commuting with
$G_3$ (Theorem~17.21 of \cite{ES17}), not the only bounded part ($b=R_{7,13}\pm R_{10,11}$ is bounded with $[b,G_3]\neq0$, \cref{d4:sec:bfield}); it
preserves the total $\tfrac12|Z|^2$ and $q_3$, hence $\tfrac12|q_3|$, but changes the constituents' rest energies, their sum included
($\tfrac12\to\tfrac32$), and the Lorentz length of the summed density (\cref{d1:prop:bfield}; d3-23, x3-22). The conclusion, interpretive and
drawn from fixed statements: the undivided energy and its differentiation in motion are complete before any $b$-field appears, and where the
hypothesis needs the closure to produce a rest energy the reading fails: for the momentum field it cannot, and the frame field's rest energy changes only
by an oscillation or an unbounded drift, with no value selected.
